\chapter{Introduzione}
I motori sincroni a magnete permanente e a riluttanza strovano applicazione un'ampia varietà di applicazioni, come nell'ambito industriale, domestico e  nell'automotive \cite{bojoi2024}.\\
I motori IPM presentano diversi vantaggi: ampia variazione di velocità e coppia, elevata potenza, basso peso ed efficienza energetica \cite{Hwang}.\\
La determinazione del fluso magnetico è fondamentale per lo sviluppo del controllo e la calibrazione \cite{bojoi2024}.\\
Spesso per la creazione del modello magnetico del motore vengono utilizzati software come FEM, ma questi ultimi necessitano una conoscenza approfondita della struttura del motore. \\
La tecnica proposta in questa tesi permette di mappare il flusso magnetico utilizzando unicamente tramite delle misurazioni eseguite sulla corrente, tensione e posizione.\\
I vantaggi dell'identificazione sperimentale del modello magnetico sono \cite{Armando}:
\begin{enumerate}
\item La performance del motore può essere calcolata offline con precisione.
\item In particolare l'utilizzo delle tecniche di controllo MTPA e MTPV sono obbligatorie per ottimizzare lo sfruttamento della coppia.
\item È possibile confrontare la performance di motori provenienti da produttori differenti, senza il bisogno di conoscere il design, materiali e fabbricazione.
\end{enumerate}


 

